\begin{thebibliography}{99}
\bibitem[1]{ref1} Michael Alekhnovich and Alexander A. Razborov. Lower bounds for polynomial calculus: Non-binomial case. In 42nd Annual Symposium on Foundations of Computer Science, FOCS 2001, 14-17 October 2001, Las Vegas, Nevada, USA, pages 190–199. IEEE Computer Society, 2001. doi:10.1109/SFCS.2001.959893.
\bibitem[2]{ref2} Yaroslav Alekseev. A lower bound for polynomial calculus with extension rule. In Valentine Kabanets, editor, 36th Computational Complexity Conference, CCC 2021, July 20-23, 2021, Toronto, Ontario, Canada (Virtual Conference), volume 200 of LIPIcs, pages 21:1–21:18. Schloss Dagstuhl – Leibniz-Zentrum für Informatik, 2021. doi:10.4230/LIPIcs.CCC.2021.21.
\bibitem[3]{ref3} Robert Andrews and Michael A. Forbes. Ideals, determinants, and straightening: Proving and using lower bounds for polynomial ideals. CoRR, abs/2112.00792, 2021. arXiv:2112.00792.
\bibitem[4]{ref4} Paul Beame, Russell Impagliazzo, Jan Krajíček, Toniann Pitassi, and Pavel Pudlák. Lower bounds on hilbert’s nullstellensatz and propositional proofs. Proceedings of the London Mathematical Society, 3(1):1–26, 1996.
\bibitem[5]{ref5} Paul Beame, Richard Karp, Toniann Pitassi, and Michael Saks. The efficiency of resolution and davis–putnam procedures. SIAM Journal on Computing, 31(4):1048–1075, 2002.
\bibitem[6]{ref6} Eli Ben-Sasson and Russell Impagliazzo. Random cnfs are hard for the polynomial calculus. In 40th Annual Symposium on Foundations of Computer Science (Cat. No. 99CB37039), pages 415–421. IEEE, 1999.
\bibitem[7]{ref7} Eli Ben-Sasson and Avi Wigderson. Short proofs are narrow – Resolution made simple. In Proceedings of the thirty-first annual ACM symposium on Theory of computing, pages 517–526, 1999.
\bibitem[8]{ref8} Sam Buss, Dima Grigoriev, Russell Impagliazzo, and Toniann Pitassi. Linear gaps between degrees for the polynomial calculus modulo distinct primes. Journal of Computer and System Sciences, 62(2):267–289, 2001.
\bibitem[9]{ref9} Vašek Chvátal and Endre Szemerédi. Many hard examples for resolution. Journal of the ACM (JACM), 35(4):759–768, 1988.
\bibitem[10]{ref10} Matthew Clegg, Jeffery Edmonds, and Russell Impagliazzo. Using the groebner basis algorithm to find proofs of unsatisfiability. In Proceedings of the twenty-eighth annual ACM symposium on Theory of computing, pages 174–183, 1996.
\bibitem[11]{ref11} Stephen A Cook and Robert A Reckhow. The relative efficiency of propositional proof systems. The journal of symbolic logic, 44(1):36–50, 1979. 7:23
\bibitem[12]{ref12} Stefan S. Dantchev and Søren Riis. On relativisation and complexity gap. In Matthias Baaz and Johann A. Makowsky, editors, Computer Science Logic, 17th International Workshop, CSL 2003, 12th Annual Conference of the EACSL, and 8th Kurt Gödel Colloquium, KGC 2003, Vienna, Austria, August 25-30, 2003, Proceedings, volume 2803 of Lecture Notes in Computer Science, pages 142–154. Springer, 2003. doi:10.1007/978-3-540-45220-1\_14.
\bibitem[13]{ref13} Michael A. Forbes, Amir Shpilka, Iddo Tzameret, and Avi Wigderson. Proof complexity lower bounds from algebraic circuit complexity. Theory Comput., 17:1–88, 2021. URL: https://theoryofcomputing.org/articles/v017a010/.
\bibitem[14]{ref14} Nashlen Govindasamy, Tuomas Hakoniemi, and Iddo Tzameret. Simple hard instances for low-depth algebraic proofs. In 63rd IEEE Annual Symposium on Foundations of Computer Science, FOCS 2022, Denver, CO, USA, October 31 – November 3, 2022, pages 188–199. IEEE, 2022.
\bibitem[15]{ref15} Armin Haken. The intractability of resolution. Theoretical computer science, 39:297–308, 1985.
\bibitem[16]{ref16} Russell Impagliazzo, Sasank Mouli, and Toniann Pitassi. The surprising power of constant depth algebraic proofs. In Proceedings of the 35th Annual ACM/IEEE Symposium on Logic in Computer Science, pages 591–603, 2020.
\bibitem[17]{ref17} Russell Impagliazzo, Pavel Pudlák, and Jirí Sgall. Lower bounds for the polynomial calculus and the gröbner basis algorithm. Comput. Complex., 8(2):127–144, 1999. doi:10.1007/ s000370050024.
\bibitem[18]{ref18} Alexander A Razborov. Lower bounds on the size of bounded depth circuits over a complete basis with logical addition. Math. notes of the Academy of Sciences of the USSR, 41(4):333–338, 1987.
\bibitem[19]{ref19} Roman Smolensky. Algebraic methods in the theory of lower bounds for boolean circuit complexity. In Proceedings of the nineteenth annual ACM symposium on Theory of computing, pages 77–82, 1987.
\bibitem[20]{ref20} Dmitry Sokolov. (semi) algebraic proofs over \{±1\} variables. In Proceedings of the 52nd Annual ACM SIGACT Symposium on Theory of Computing, pages 78–90, 2020.
\end{thebibliography}
